\documentclass[12pt]{article}
\usepackage{amsmath,amssymb,amsfonts}
\usepackage[margin=1in]{geometry}

\begin{document}

\textbf{Chapter 8, Problem 23.} Show that, by an obvious definition of addition and scalar multiplication, the bounded linear functionals of $L_p$ form a Banach space in their own dual space. This space is called the \emph{dual space} of $B$ and is denoted by $B^*$. Since every linear functional on $L_p$ is generated by an element of $L_q$ ($1/p+1/q=1$), we can identify $L_p^*$ and $L_q$, and say that $L_q$ is the dual space of $L_p$.

Note that the space of square-integrable functions is its own dual space.

(There are three parts to this problem. First, you must show that your definition of addition and scalar multiplication is such that it gives rise to a linear vector space. Second, you must show that the definition of norm given in the previous problem really is a norm (as we did in Section 8.2 in discussing norms of operators). There parts are easy. No so easy is the problem of showing that the normed linear vector space which you have constructed is \emph{complete}.)

\bigskip
\textbf{Solution.}

\textbf{Part 1: $B^*$ is a linear vector space.}

Define addition and scalar multiplication of bounded linear functionals on $B$:
\begin{itemize}
\item $(F_1+F_2)(g) = F_1(g)+F_2(g)$ for all $g\in B$.
\item $(\alpha F)(g) = \alpha F(g)$ for all $g\in B$, $\alpha\in\mathbb{C}$.
\end{itemize}

These operations are well-defined and satisfy the vector space axioms:
\begin{itemize}
\item $F_1+F_2$ is linear: $(F_1+F_2)(\alpha g+\beta h) = \alpha(F_1+F_2)(g)+\beta(F_1+F_2)(h)$.
\item $F_1+F_2$ is bounded: $|(F_1+F_2)(g)|\le|F_1(g)|+|F_2(g)|\le(\|F_1\|+\|F_2\|)\|g\|$.
\item $\alpha F$ is linear and bounded with $\|\alpha F\| = |\alpha|\|F\|$.
\item The zero functional $F\equiv 0$ serves as the additive identity.
\end{itemize}
All other axioms (associativity, commutativity, etc.) follow trivially. $\checkmark$

\textbf{Part 2: $\|F\|$ is a norm.}

We verify the norm axioms:
\begin{enumerate}
\item $\|F\|\ge 0$ and $\|F\|=0\iff F=0$: Clear from the definition $\|F\| = \sup_{\|g\|=1}|F(g)|$.
\item $\|\alpha F\| = |\alpha|\|F\|$: $\|\alpha F\| = \sup|(\alpha F)(g)|/\|g\| = |\alpha|\sup|F(g)|/\|g\| = |\alpha|\|F\|$.
\item Triangle inequality: $\|F_1+F_2\| = \sup\frac{|(F_1+F_2)(g)|}{\|g\|} \le \sup\frac{|F_1(g)|+|F_2(g)|}{\|g\|}\le \|F_1\|+\|F_2\|$. $\checkmark$
\end{enumerate}

\textbf{Part 3: $B^*$ is complete.}

Let $\{F_n\}$ be a Cauchy sequence in $B^*$, i.e., $\|F_n-F_m\|\to 0$ as $n,m\to\infty$.

\emph{Step 1:} For each fixed $g\in B$:
\[
|F_n(g)-F_m(g)| = |(F_n-F_m)(g)| \le \|F_n-F_m\|\|g\| \to 0.
\]
So $\{F_n(g)\}$ is a Cauchy sequence of complex numbers and converges. Define $F(g) = \lim_{n\to\infty}F_n(g)$.

\emph{Step 2:} $F$ is linear: $F(\alpha g+\beta h) = \lim F_n(\alpha g+\beta h) = \alpha\lim F_n(g)+\beta\lim F_n(h) = \alpha F(g)+\beta F(h)$.

\emph{Step 3:} $F$ is bounded: Since $\{F_n\}$ is Cauchy, $\{\|F_n\|\}$ is bounded, say $\|F_n\|\le M$ for all $n$. Then $|F(g)| = \lim|F_n(g)|\le M\|g\|$.

\emph{Step 4:} $F_n\to F$ in norm: Given $\varepsilon>0$, choose $N$ such that $\|F_n-F_m\|<\varepsilon$ for $n,m>N$. For any $g$ with $\|g\|=1$:
\[
|F_n(g)-F_m(g)| < \varepsilon.
\]
Taking $m\to\infty$: $|F_n(g)-F(g)|\le\varepsilon$ for all $\|g\|=1$ and $n>N$. Thus $\|F_n-F\|\le\varepsilon$.

Therefore $B^*$ is complete, hence a Banach space.

\textbf{Identification:} Since every bounded linear functional on $L_p$ is of the form $F(g)=\int fg\,dx$ for a unique $f\in L_q$ (by the Riesz representation theorem), and $\|F\|=\|f\|_q$ (Problem 22), we have $L_p^* \cong L_q$ isometrically.

For $p=q=2$: $L_2^* \cong L_2$, so Hilbert space is its own dual (this is the Riesz representation theorem for Hilbert spaces). $\blacksquare$

\end{document}
